\subsubsection{突発型 EEJ の発生要因}
\label{sec:discussion_brazil_sudden}

突発型 EEJ は、未発達型と比較して季節依存性が弱いものの、
9 月および 10 月において発生率が約 7\% と高い傾向が確認された。
一方で、D-solstice では発生率が低い結果となった。

この結果は、Sq 電流の季節変動が突発型 EEJ の発生に一定の影響を及ぼしているものの、
Sq 電流の強弱によって、
Dip 観測点のEUELを大きく押し下げる直接的な要因ではないことを示唆している。

一方、月潮汐効果においては、
Lunar age = 6--7 および 17--20 付近で発生率が高くなる傾向が見られた。
これは、月潮汐効果によって Dip 観測点における西向き電流が卓越し、
EUEL が大きく減少した結果、
Off-dip 観測点と比較して顕著な差が生じたためであると解釈できる。

以上の結果から、突発型 EEJ は、
J-solstice に代表される Sq 電流が相対的に弱い時期において、
月潮汐効果による Dip 観測点での西向き電流の増強が重畳することで発生すると考えられる。
